\documentclass[UTF8,11pt]{ctexart}

\usepackage[a4paper,margin=2.2cm]{geometry}
\usepackage{amsmath,amssymb,bm}
\usepackage{booktabs}
\usepackage{array}
\usepackage{hyperref}
\usepackage{xcolor}
\usepackage{listings}

\lstset{
  basicstyle=\ttfamily\small,
  breaklines=true,
  frame=single,
  columns=fullflexible
}

\title{SPHinXsys 中 A--$\phi$ Matrix-free 求解器的粒子离散与迭代求解说明}
\author{}
\date{}

\begin{document}
\maketitle

\section{先明确一个基本关系：Residual = RHS - LHS}

在离散线性系统中，方程可以写成
\begin{equation}
    \mathcal{K} X = B .
\end{equation}
其中：
\begin{itemize}
    \item $X$ 是未知量，也就是当前要求解的粒子场；
    \item $B$ 是右端项，即 RHS，通常来自外加源项、电流源、边界项等；
    \item $\mathcal{K}$ 是离散后的 A--$\phi$ 线性算子；
    \item $\mathcal{K}X$ 是把当前未知量 $X$ 代入方程左端以后得到的值，即 LHS。
\end{itemize}

因此残差定义为
\begin{equation}
    R = B - \mathcal{K}X .
\end{equation}

也就是说，程序里的
\begin{lstlisting}
Residual = RHS - LHS;
\end{lstlisting}
就是
\begin{equation}
    \mathrm{Residual} = \mathrm{RightHandSide} - \mathrm{LeftHandSide}.
\end{equation}

如果当前解已经满足离散方程，则
\begin{equation}
    \mathcal{K}X = B,
\end{equation}
所以
\begin{equation}
    R = 0 .
\end{equation}

这里特别注意：LHS 不是残差。LHS 是当前解代入方程左端得到的值，Residual 才是 RHS 与 LHS 之间的差值。

\section{连续 A--$\phi$ 方程与实虚部分裂}

这里采用频域 A--$\phi$ 方程的一种常用形式：
\begin{equation}
    -\nabla\cdot(\nu\nabla \bm A)
    + j\omega\sigma \bm A
    + \sigma\nabla \phi
    = \bm J_s ,
    \label{eq:continuous-a}
\end{equation}
\begin{equation}
    -\nabla\cdot(\sigma\nabla \phi)
    - j\omega\nabla\cdot(\sigma \bm A)
    = 0 .
    \label{eq:continuous-phi}
\end{equation}

其中：
\begin{itemize}
    \item $\bm A$ 是磁矢势；
    \item $\phi$ 是电势；
    \item $\nu = 1/\mu$ 是磁阻率；
    \item $\sigma$ 是电导率；
    \item $\omega$ 是角频率；
    \item $\bm J_s$ 是外加源电流密度；
    \item $j=\sqrt{-1}$。
\end{itemize}

由于 SPHinXsys 当前的粒子变量系统更自然地支持实数场和向量场，因此可以把复数变量拆成实部和虚部：
\begin{equation}
    \bm A = \bm A_{\mathrm{re}} + j\bm A_{\mathrm{im}},
\end{equation}
\begin{equation}
    \phi = \phi_{\mathrm{re}} + j\phi_{\mathrm{im}},
\end{equation}
\begin{equation}
    \bm J_s = \bm J_{s,\mathrm{re}} + j\bm J_{s,\mathrm{im}}.
\end{equation}

将
\begin{equation}
    j\omega\sigma\bm A
    =
    j\omega\sigma
    \left(
        \bm A_{\mathrm{re}} + j\bm A_{\mathrm{im}}
    \right)
    =
    -\omega\sigma\bm A_{\mathrm{im}}
    +
    j\omega\sigma\bm A_{\mathrm{re}}
\end{equation}
代入方程，可得到四个实数方程。

A 方程实部：
\begin{equation}
    -\nabla\cdot(\nu\nabla \bm A_{\mathrm{re}})
    - \omega\sigma\bm A_{\mathrm{im}}
    + \sigma\nabla\phi_{\mathrm{re}}
    =
    \bm J_{s,\mathrm{re}} .
    \label{eq:a-real}
\end{equation}

A 方程虚部：
\begin{equation}
    -\nabla\cdot(\nu\nabla \bm A_{\mathrm{im}})
    + \omega\sigma\bm A_{\mathrm{re}}
    + \sigma\nabla\phi_{\mathrm{im}}
    =
    \bm J_{s,\mathrm{im}} .
    \label{eq:a-imag}
\end{equation}

$\phi$ 方程实部：
\begin{equation}
    -\nabla\cdot(\sigma\nabla \phi_{\mathrm{re}})
    + \omega\nabla\cdot(\sigma\bm A_{\mathrm{im}})
    =
    0 .
    \label{eq:phi-real}
\end{equation}

$\phi$ 方程虚部：
\begin{equation}
    -\nabla\cdot(\sigma\nabla \phi_{\mathrm{im}})
    - \omega\nabla\cdot(\sigma\bm A_{\mathrm{re}})
    =
    0 .
    \label{eq:phi-imag}
\end{equation}

\section{粒子上的未知量组织}

假设共有 $N$ 个粒子。对三维问题，每个粒子 $i$ 上的未知量为
\begin{equation}
    \bm A_{\mathrm{re},i}
    =
    \begin{bmatrix}
    A^x_{\mathrm{re},i}\\
    A^y_{\mathrm{re},i}\\
    A^z_{\mathrm{re},i}
    \end{bmatrix},
    \qquad
    \bm A_{\mathrm{im},i}
    =
    \begin{bmatrix}
    A^x_{\mathrm{im},i}\\
    A^y_{\mathrm{im},i}\\
    A^z_{\mathrm{im},i}
    \end{bmatrix},
\end{equation}
以及两个标量
\begin{equation}
    \phi_{\mathrm{re},i},
    \qquad
    \phi_{\mathrm{im},i}.
\end{equation}

因此每个粒子上有 $8$ 个自由度：
\begin{equation}
    X_i
    =
    \begin{bmatrix}
    A^x_{\mathrm{re},i}\\
    A^y_{\mathrm{re},i}\\
    A^z_{\mathrm{re},i}\\
    A^x_{\mathrm{im},i}\\
    A^y_{\mathrm{im},i}\\
    A^z_{\mathrm{im},i}\\
    \phi_{\mathrm{re},i}\\
    \phi_{\mathrm{im},i}
    \end{bmatrix}.
\end{equation}

所有粒子的未知量合起来是
\begin{equation}
    X =
    \begin{bmatrix}
    X_1\\
    X_2\\
    \vdots\\
    X_N
    \end{bmatrix}.
\end{equation}

所以如果显式组装全局矩阵，那么矩阵规模大约是
\begin{equation}
    8N \times 8N.
\end{equation}

但是 matrix-free 方法不显式存储这个 $8N\times 8N$ 的全局稀疏矩阵，而是只实现一个粒子邻居计算过程，用来计算
\begin{equation}
    Y = \mathcal{K}X .
\end{equation}

\section{SPH 粒子几何与 kernel 记号}

令粒子位置为 $\bm r_i$，邻居粒子为 $j\in\mathcal{N}(i)$。定义
\begin{equation}
    \bm r_{ij} = \bm r_i - \bm r_j,
    \qquad
    r_{ij} = |\bm r_{ij}|,
    \qquad
    \bm e_{ij} = \frac{\bm r_{ij}}{r_{ij}} .
\end{equation}

令
\begin{equation}
    W_{ij} = W(r_{ij},h),
    \qquad
    dW_{ij} = \frac{\partial W(r_{ij},h)}{\partial r_{ij}} .
\end{equation}

SPHinXsys 中还会使用：
\begin{itemize}
    \item $V_j$：邻居粒子体积；
    \item $\bm B_i$：粒子 $i$ 上的 kernel 修正矩阵；
    \item $\bm B_j$：粒子 $j$ 上的 kernel 修正矩阵。
\end{itemize}

为了表达简洁，定义一个梯度型权重向量：
\begin{equation}
    \bm G_{ij}
    =
    V_j \bm B_i \bm e_{ij} dW_{ij}.
    \label{eq:gij}
\end{equation}

在 SPHinXsys 的实际代码中，$\bm e_{ij}$、$dW_{ij}$ 的符号约定要与库内部保持一致。本文档中的公式主要说明结构；实现时应以现有 \texttt{LinearGradient} 和 \texttt{diffusion\_dynamics\_ck} 的符号约定为准。

对材料系数，建议即便单材料也按粒子数组存储：
\begin{equation}
    \sigma_i,
    \qquad
    \nu_i.
\end{equation}

多材料界面上可以使用调和平均，例如
\begin{equation}
    \nu_{ij}
    =
    \frac{2\nu_i\nu_j}{\nu_i+\nu_j+\epsilon},
\end{equation}
\begin{equation}
    \sigma_{ij}
    =
    \frac{2\sigma_i\sigma_j}{\sigma_i+\sigma_j+\epsilon}.
\end{equation}

\section{粒子形式的梯度、散度与 Laplace 型算子}

\subsection{梯度}

对标量场 $\psi$，可以写成
\begin{equation}
    (\nabla\psi)_i
    \approx
    \sum_{j\in\mathcal{N}(i)}
    \bm G_{ij}
    \left(
        \psi_i - \psi_j
    \right).
    \label{eq:grad-scalar}
\end{equation}

因此
\begin{equation}
    (\nabla\phi_{\mathrm{re}})_i
    \approx
    \sum_{j\in\mathcal{N}(i)}
    \bm G_{ij}
    \left(
        \phi_{\mathrm{re},i}
        -
        \phi_{\mathrm{re},j}
    \right),
\end{equation}
\begin{equation}
    (\nabla\phi_{\mathrm{im}})_i
    \approx
    \sum_{j\in\mathcal{N}(i)}
    \bm G_{ij}
    \left(
        \phi_{\mathrm{im},i}
        -
        \phi_{\mathrm{im},j}
    \right).
\end{equation}

\subsection{散度}

对向量场 $\bm q$，可以写成
\begin{equation}
    (\nabla\cdot\bm q)_i
    \approx
    \sum_{j\in\mathcal{N}(i)}
    \bm G_{ij}\cdot
    \left(
        \bm q_i - \bm q_j
    \right).
    \label{eq:div-vector}
\end{equation}

对于 $\nabla\cdot(\sigma\bm A)$，可以把电导率放进 pairwise 系数中：
\begin{equation}
    \left[\nabla\cdot(\sigma\bm A)\right]_i
    \approx
    \sum_{j\in\mathcal{N}(i)}
    \sigma_{ij}\bm G_{ij}\cdot
    \left(
        \bm A_i - \bm A_j
    \right).
    \label{eq:div-sigma-a}
\end{equation}

因此
\begin{equation}
    \left[\nabla\cdot(\sigma\bm A_{\mathrm{re}})\right]_i
    \approx
    \sum_{j\in\mathcal{N}(i)}
    \sigma_{ij}\bm G_{ij}\cdot
    \left(
        \bm A_{\mathrm{re},i}
        -
        \bm A_{\mathrm{re},j}
    \right),
\end{equation}
\begin{equation}
    \left[\nabla\cdot(\sigma\bm A_{\mathrm{im}})\right]_i
    \approx
    \sum_{j\in\mathcal{N}(i)}
    \sigma_{ij}\bm G_{ij}\cdot
    \left(
        \bm A_{\mathrm{im},i}
        -
        \bm A_{\mathrm{im},j}
    \right).
\end{equation}

\subsection{Laplace 型算子}

A--$\phi$ 方程中需要的主算子是
\begin{equation}
    -\nabla\cdot(\kappa\nabla u),
\end{equation}
其中 $\kappa$ 可以是 $\nu$ 或 $\sigma$。

为了与 matrix-free 求解器配合，建议把它写成 pairwise 形式：
\begin{equation}
    \left(\mathcal{L}_{\kappa}u\right)_i
    =
    \left[-\nabla\cdot(\kappa\nabla u)\right]_i
    \approx
    \sum_{j\in\mathcal{N}(i)}
    c_{ij}^{\kappa}
    \left(
        u_i - u_j
    \right).
    \label{eq:laplace-pairwise}
\end{equation}

其中 $c_{ij}^{\kappa}$ 是 Laplace/diffusion pairwise 权重。一个常见的 SPH 扩散形式可以抽象写成
\begin{equation}
    c_{ij}^{\kappa}
    =
    \kappa_{ij}
    \frac{\bm r_{ij}}{r_{ij}^2+\eta h^2}
    \cdot
    \bm S_{ij},
    \label{eq:laplace-weight-general}
\end{equation}
其中 $\eta$ 是小的正则化常数，$\bm S_{ij}$ 是与 kernel 梯度和粒子体积有关的面积型向量，例如可采用与 SPHinXsys diffusion 算子一致的形式：
\begin{equation}
    \bm S_{ij}
    =
    V_j
    \left(\bm B_i+\bm B_j\right)
    \bm e_{ij}dW_{ij}.
    \label{eq:surface-area-like}
\end{equation}

实现中不一定要逐字使用上式，而是应该直接参考 SPHinXsys 中 \texttt{diffusion\_dynamics\_ck} 的 pairwise diffusion 公式。关键是最终形式应当保持为
\begin{equation}
    \left(\mathcal{L}_{\kappa}u\right)_i
    =
    \sum_j c_{ij}^{\kappa}(u_i-u_j),
\end{equation}
这样每个线程只计算并写入粒子 $i$ 的结果。

对向量场 $\bm A$，可以对每个分量使用同一个 Laplace 型算子：
\begin{equation}
    \left(\mathcal{L}_{\nu}\bm A\right)_i^m
    =
    \sum_{j\in\mathcal{N}(i)}
    c_{ij}^{\nu}
    \left(
        A_i^m-A_j^m
    \right),
    \qquad
    m=x,y,z.
\end{equation}

\section{A--$\phi$ 左端项的粒子离散形式}

定义
\begin{equation}
    \mathcal{L}_{\nu}\bm A
    \approx
    -\nabla\cdot(\nu\nabla \bm A),
\end{equation}
\begin{equation}
    \mathcal{L}_{\sigma}\phi
    \approx
    -\nabla\cdot(\sigma\nabla \phi).
\end{equation}

对 $A$ 的第 $m$ 个分量，$m=x,y,z$，A 方程实部的左端项为
\begin{align}
    \mathrm{LhsAReal}_i^m
    &=
    \sum_{j\in\mathcal{N}(i)}
    c_{ij}^{\nu}
    \left(
        A_{\mathrm{re},i}^m
        -
        A_{\mathrm{re},j}^m
    \right)
    -
    \omega\sigma_i A_{\mathrm{im},i}^m
    \nonumber\\
    &\quad
    +
    \sigma_i
    \sum_{j\in\mathcal{N}(i)}
    G_{ij}^m
    \left(
        \phi_{\mathrm{re},i}
        -
        \phi_{\mathrm{re},j}
    \right).
    \label{eq:lhs-a-real-particle}
\end{align}

A 方程虚部的左端项为
\begin{align}
    \mathrm{LhsAImag}_i^m
    &=
    \sum_{j\in\mathcal{N}(i)}
    c_{ij}^{\nu}
    \left(
        A_{\mathrm{im},i}^m
        -
        A_{\mathrm{im},j}^m
    \right)
    +
    \omega\sigma_i A_{\mathrm{re},i}^m
    \nonumber\\
    &\quad
    +
    \sigma_i
    \sum_{j\in\mathcal{N}(i)}
    G_{ij}^m
    \left(
        \phi_{\mathrm{im},i}
        -
        \phi_{\mathrm{im},j}
    \right).
    \label{eq:lhs-a-imag-particle}
\end{align}

$\phi$ 方程实部的左端项为
\begin{align}
    \mathrm{LhsPhiReal}_i
    &=
    \sum_{j\in\mathcal{N}(i)}
    c_{ij}^{\sigma}
    \left(
        \phi_{\mathrm{re},i}
        -
        \phi_{\mathrm{re},j}
    \right)
    \nonumber\\
    &\quad
    +
    \omega
    \sum_{j\in\mathcal{N}(i)}
    \sigma_{ij}\bm G_{ij}\cdot
    \left(
        \bm A_{\mathrm{im},i}
        -
        \bm A_{\mathrm{im},j}
    \right).
    \label{eq:lhs-phi-real-particle}
\end{align}

$\phi$ 方程虚部的左端项为
\begin{align}
    \mathrm{LhsPhiImag}_i
    &=
    \sum_{j\in\mathcal{N}(i)}
    c_{ij}^{\sigma}
    \left(
        \phi_{\mathrm{im},i}
        -
        \phi_{\mathrm{im},j}
    \right)
    \nonumber\\
    &\quad
    -
    \omega
    \sum_{j\in\mathcal{N}(i)}
    \sigma_{ij}\bm G_{ij}\cdot
    \left(
        \bm A_{\mathrm{re},i}
        -
        \bm A_{\mathrm{re},j}
    \right).
    \label{eq:lhs-phi-imag-particle}
\end{align}

以上四组量就是程序中 \texttt{AphiApply} 要计算的内容。也就是
\begin{equation}
    \mathrm{LHS} = \mathcal{K}(X).
\end{equation}

\section{RHS 与残差的粒子形式}

对粒子 $i$，右端项可以写成
\begin{equation}
    B_i
    =
    \begin{bmatrix}
    J^x_{s,\mathrm{re},i}\\
    J^y_{s,\mathrm{re},i}\\
    J^z_{s,\mathrm{re},i}\\
    J^x_{s,\mathrm{im},i}\\
    J^y_{s,\mathrm{im},i}\\
    J^z_{s,\mathrm{im},i}\\
    B_{\phi,\mathrm{re},i}\\
    B_{\phi,\mathrm{im},i}
    \end{bmatrix}.
\end{equation}

一般情况下，$\phi$ 方程没有体源项，因此
\begin{equation}
    B_{\phi,\mathrm{re},i}=0,
    \qquad
    B_{\phi,\mathrm{im},i}=0.
\end{equation}

如果边界条件、约束项或特殊源项会进入 $\phi$ 方程右端，则这两个量可以不是零。

左端项为
\begin{equation}
    (\mathcal{K}X)_i
    =
    \begin{bmatrix}
    \mathrm{LhsAReal}_i^x\\
    \mathrm{LhsAReal}_i^y\\
    \mathrm{LhsAReal}_i^z\\
    \mathrm{LhsAImag}_i^x\\
    \mathrm{LhsAImag}_i^y\\
    \mathrm{LhsAImag}_i^z\\
    \mathrm{LhsPhiReal}_i\\
    \mathrm{LhsPhiImag}_i
    \end{bmatrix}.
\end{equation}

残差为
\begin{equation}
    R_i
    =
    B_i
    -
    (\mathcal{K}X)_i .
\end{equation}

展开为
\begin{equation}
    R_i
    =
    \begin{bmatrix}
    J^x_{s,\mathrm{re},i} - \mathrm{LhsAReal}_i^x\\
    J^y_{s,\mathrm{re},i} - \mathrm{LhsAReal}_i^y\\
    J^z_{s,\mathrm{re},i} - \mathrm{LhsAReal}_i^z\\
    J^x_{s,\mathrm{im},i} - \mathrm{LhsAImag}_i^x\\
    J^y_{s,\mathrm{im},i} - \mathrm{LhsAImag}_i^y\\
    J^z_{s,\mathrm{im},i} - \mathrm{LhsAImag}_i^z\\
    B_{\phi,\mathrm{re},i} - \mathrm{LhsPhiReal}_i\\
    B_{\phi,\mathrm{im},i} - \mathrm{LhsPhiImag}_i
    \end{bmatrix}.
\end{equation}

程序中对应：
\begin{lstlisting}
ResidualAReal[i]   = RhsAReal[i]   - LhsAReal[i];
ResidualAImag[i]   = RhsAImag[i]   - LhsAImag[i];
ResidualPhiReal[i] = RhsPhiReal[i] - LhsPhiReal[i];
ResidualPhiImag[i] = RhsPhiImag[i] - LhsPhiImag[i];
\end{lstlisting}

\section{如果显式组装矩阵，它大概长什么样？}

如果把所有粒子的未知量都排成全局向量，可以抽象写成
\begin{equation}
\begin{bmatrix}
L_{\nu} & -\omega\Sigma & \Sigma G & 0\\
\omega\Sigma & L_{\nu} & 0 & \Sigma G\\
0 & \omega D_{\sigma} & L_{\sigma} & 0\\
-\omega D_{\sigma} & 0 & 0 & L_{\sigma}
\end{bmatrix}
\begin{bmatrix}
A_{\mathrm{re}}\\
A_{\mathrm{im}}\\
\phi_{\mathrm{re}}\\
\phi_{\mathrm{im}}
\end{bmatrix}
=
\begin{bmatrix}
J_{s,\mathrm{re}}\\
J_{s,\mathrm{im}}\\
B_{\phi,\mathrm{re}}\\
B_{\phi,\mathrm{im}}
\end{bmatrix}.
\label{eq:block-global}
\end{equation}

这里：
\begin{itemize}
    \item $L_{\nu}$ 表示 $-\nabla\cdot(\nu\nabla \bm A)$ 对应的全局 Laplace 型稀疏矩阵；
    \item $L_{\sigma}$ 表示 $-\nabla\cdot(\sigma\nabla\phi)$ 对应的全局 Laplace 型稀疏矩阵；
    \item $\Sigma$ 是由 $\sigma_i$ 组成的对角矩阵；
    \item $G$ 是 gradient 算子；
    \item $D_{\sigma}$ 是 $\nabla\cdot(\sigma\bm A)$ 对应的 divergence 算子。
\end{itemize}

这里需要特别注意：$\phi$ 方程不是没有系数矩阵。$\phi$ 自身对应的系数矩阵是 $L_{\sigma}$。在式 \eqref{eq:block-global} 中，第三行第三列和第四行第四列的 $L_{\sigma}$ 就是 $\phi_{\mathrm{re}}$ 和 $\phi_{\mathrm{im}}$ 的主系数矩阵。

$D_{\sigma}$ 只是 A 场进入 $\phi$ 方程的耦合项：
\begin{equation}
    \omega D_{\sigma}A_{\mathrm{im}},
    \qquad
    -\omega D_{\sigma}A_{\mathrm{re}}.
\end{equation}

因此 $\phi$ 方程的结构是
\begin{equation}
    L_{\sigma}\phi_{\mathrm{re}}
    +
    \omega D_{\sigma}A_{\mathrm{im}}
    =
    B_{\phi,\mathrm{re}},
\end{equation}
\begin{equation}
    L_{\sigma}\phi_{\mathrm{im}}
    -
    \omega D_{\sigma}A_{\mathrm{re}}
    =
    B_{\phi,\mathrm{im}}.
\end{equation}

\section{不生成稀疏矩阵，为什么还能求解？}

Matrix-free 的核心思想是：求解器不需要显式知道矩阵中每一个元素，只需要能计算
\begin{equation}
    Y = \mathcal{K}X .
\end{equation}

如果显式矩阵存在，那么第 $i$ 行的矩阵乘法是
\begin{equation}
    Y_i
    =
    \sum_j K_{ij}X_j .
\end{equation}

在 SPH 中，$K_{ij}$ 的非零项只来自粒子 $i$ 的邻居 $j\in\mathcal{N}(i)$。所以可以不存 $K_{ij}$，而是在每次需要 $Y=\mathcal{K}X$ 的时候，通过邻居循环直接计算：
\begin{equation}
    Y_i
    =
    \sum_{j\in\mathcal{N}(i)}
    \text{pairwise contribution}(i,j;X_i,X_j).
\end{equation}

程序结构就是：
\begin{lstlisting}
parallel_for each particle i:
    lhs_i = 0

    for each neighbor j of i:
        lhs_i += pairwise_laplace_terms(i, j)
        lhs_i += pairwise_gradient_coupling(i, j)
        lhs_i += pairwise_divergence_coupling(i, j)

    lhs_i += local_reaction_terms(i)

    LHS[i] = lhs_i
\end{lstlisting}

所以不是“不求矩阵乘法”，而是“不存矩阵，直接按粒子邻居公式计算矩阵乘法的结果”。

Krylov / BiCGStab 等迭代方法只需要反复调用
\begin{equation}
    Y = \mathcal{K}X
\end{equation}
和做若干向量加法、内积、范数计算，因此非常适合 matrix-free 实现。

\section{Krylov 搜索方向与 BiCGStab 中的变量含义}

对于完整 A--$\phi$ 系统，由于实虚部耦合、A--$\phi$ 耦合和 gauge 约束，整体算子通常不是对称正定的，因此不建议直接使用普通 CG。更合适的第一选择是 BiCGStab。

令
\begin{equation}
    X^k
\end{equation}
表示第 $k$ 次线性求解迭代时的当前解。它包含四组粒子场：
\begin{equation}
    X^k =
    \left[
    \bm A_{\mathrm{re}}^k,
    \bm A_{\mathrm{im}}^k,
    \phi_{\mathrm{re}}^k,
    \phi_{\mathrm{im}}^k
    \right]^T .
\end{equation}

残差为
\begin{equation}
    R^k = B - \mathcal{K}X^k .
\end{equation}

搜索方向 $P^k$ 可以理解为：当前迭代准备沿着哪个方向去修正解。初始时通常可以取
\begin{equation}
    P^0 = R^0 .
\end{equation}

但从第二步开始，搜索方向不再简单等于残差，而是结合当前残差和之前的搜索方向。BiCGStab 中有
\begin{equation}
    P^k
    =
    R^k
    +
    \beta_k
    \left(
        P^{k-1}
        -
        \omega_{k-1}V^{k-1}
    \right),
\end{equation}
其中
\begin{equation}
    V^{k-1}
    =
    \mathcal{K}P^{k-1}.
\end{equation}

也就是说，$V$ 不是速度，也不是物理场，而是搜索方向 $P$ 经过 A--$\phi$ 线性算子作用后的结果。

BiCGStab 的一个常用流程如下：

\begin{align}
    R^0 &= B-\mathcal{K}X^0,\\
    \widehat{R} &= R^0,\\
    \rho_k &= \langle \widehat{R},R^k\rangle,\\
    \beta_k &=
    \frac{\rho_k}{\rho_{k-1}}
    \frac{\alpha_{k-1}}{\omega_{k-1}},\\
    P^k &=
    R^k+\beta_k(P^{k-1}-\omega_{k-1}V^{k-1}),\\
    V^k &= \mathcal{K}P^k,\\
    \alpha_k &=
    \frac{\rho_k}{\langle \widehat{R},V^k\rangle},\\
    S^k &= R^k-\alpha_k V^k,\\
    T^k &= \mathcal{K}S^k,\\
    \omega_k &=
    \frac{\langle T^k,S^k\rangle}
         {\langle T^k,T^k\rangle},\\
    X^{k+1}
    &=
    X^k+\alpha_k P^k+\omega_k S^k,\\
    R^{k+1}
    &=
    S^k-\omega_k T^k .
\end{align}

其中：
\begin{itemize}
    \item $\alpha_k$ 是沿搜索方向 $P^k$ 的步长；
    \item $\omega_k$ 是对中间残差 $S^k$ 的进一步修正步长；
    \item $\beta_k$ 用来混合当前残差和历史搜索方向；
    \item $V^k=\mathcal{K}P^k$；
    \item $T^k=\mathcal{K}S^k$。
\end{itemize}

\section{为什么有时可以写 $R^{k+1}=R^k-\alpha\mathcal{K}P^k$？}

从残差定义出发：
\begin{equation}
    R^k = B-\mathcal{K}X^k.
\end{equation}

如果某一步只做简单更新
\begin{equation}
    X^{k+1}=X^k+\alpha P^k,
\end{equation}
那么
\begin{align}
    R^{k+1}
    &=
    B-\mathcal{K}X^{k+1}
    \nonumber\\
    &=
    B-\mathcal{K}(X^k+\alpha P^k)
    \nonumber\\
    &=
    B-\mathcal{K}X^k-\alpha\mathcal{K}P^k
    \nonumber\\
    &=
    R^k-\alpha\mathcal{K}P^k.
\end{align}

所以
\begin{equation}
    R^{k+1}=R^k-\alpha\mathcal{K}P^k
\end{equation}
并不矛盾，它是由
\begin{equation}
    R^k = B-\mathcal{K}X^k
\end{equation}
推出来的。

在 BiCGStab 中，由于还会引入 $S^k$ 和 $\omega_k$ 修正，因此最终残差更新写成
\begin{equation}
    R^{k+1}=S^k-\omega_k T^k.
\end{equation}

\section{A--$\phi$ 的粒子级 BiCGStab 操作展开}

每一个 Krylov 向量都是一个 block 粒子场。例如
\begin{equation}
    P_i^k
    =
    \begin{bmatrix}
    P^x_{A,\mathrm{re},i}\\
    P^y_{A,\mathrm{re},i}\\
    P^z_{A,\mathrm{re},i}\\
    P^x_{A,\mathrm{im},i}\\
    P^y_{A,\mathrm{im},i}\\
    P^z_{A,\mathrm{im},i}\\
    P_{\phi,\mathrm{re},i}\\
    P_{\phi,\mathrm{im},i}
    \end{bmatrix}^k.
\end{equation}

解更新
\begin{equation}
    X^{k+1}=X^k+\alpha_k P^k+\omega_k S^k
\end{equation}
在粒子 $i$ 上展开为
\begin{align}
    \bm A_{\mathrm{re},i}^{k+1}
    &=
    \bm A_{\mathrm{re},i}^{k}
    +
    \alpha_k \bm P_{A,\mathrm{re},i}^{k}
    +
    \omega_k \bm S_{A,\mathrm{re},i}^{k},
    \\
    \bm A_{\mathrm{im},i}^{k+1}
    &=
    \bm A_{\mathrm{im},i}^{k}
    +
    \alpha_k \bm P_{A,\mathrm{im},i}^{k}
    +
    \omega_k \bm S_{A,\mathrm{im},i}^{k},
    \\
    \phi_{\mathrm{re},i}^{k+1}
    &=
    \phi_{\mathrm{re},i}^{k}
    +
    \alpha_k P_{\phi,\mathrm{re},i}^{k}
    +
    \omega_k S_{\phi,\mathrm{re},i}^{k},
    \\
    \phi_{\mathrm{im},i}^{k+1}
    &=
    \phi_{\mathrm{im},i}^{k}
    +
    \alpha_k P_{\phi,\mathrm{im},i}^{k}
    +
    \omega_k S_{\phi,\mathrm{im},i}^{k}.
\end{align}

内积可以定义为粒子体积加权的 block 内积：
\begin{align}
    \langle X,Y\rangle
    =
    \sum_i V_i
    \Big[
    &\bm A_{\mathrm{re},X,i}\cdot\bm A_{\mathrm{re},Y,i}
    +
    \bm A_{\mathrm{im},X,i}\cdot\bm A_{\mathrm{im},Y,i}
    \nonumber\\
    &+
    \phi_{\mathrm{re},X,i}\phi_{\mathrm{re},Y,i}
    +
    \phi_{\mathrm{im},X,i}\phi_{\mathrm{im},Y,i}
    \Big].
\end{align}

残差范数为
\begin{equation}
    \|R\|
    =
    \sqrt{\langle R,R\rangle}.
\end{equation}

\section{程序实现上的建议}

在 SPHinXsys 的 CK/SYCL 框架中，可以把核心实现分成以下部分：

\subsection{AphiApply}

输入：
\begin{lstlisting}
AReal, AImag, PhiReal, PhiImag
\end{lstlisting}

输出：
\begin{lstlisting}
LhsAReal, LhsAImag, LhsPhiReal, LhsPhiImag
\end{lstlisting}

粒子循环结构：
\begin{lstlisting}
parallel_for each particle i:
    lhs_A_re = Vecd::Zero()
    lhs_A_im = Vecd::Zero()
    lhs_phi_re = 0
    lhs_phi_im = 0

    for each neighbor j:
        compute c_nu_ij
        compute c_sigma_ij
        compute G_ij

        lhs_A_re += Laplace(AReal)
        lhs_A_im += Laplace(AImag)
        lhs_phi_re += Laplace(PhiReal)
        lhs_phi_im += Laplace(PhiImag)

        lhs_A_re += sigma_i * grad(PhiReal)
        lhs_A_im += sigma_i * grad(PhiImag)

        lhs_phi_re += omega * div(sigma * AImag)
        lhs_phi_im -= omega * div(sigma * AReal)

    lhs_A_re += -omega * sigma_i * AImag[i]
    lhs_A_im +=  omega * sigma_i * AReal[i]

    LhsAReal[i] = lhs_A_re
    LhsAImag[i] = lhs_A_im
    LhsPhiReal[i] = lhs_phi_re
    LhsPhiImag[i] = lhs_phi_im
\end{lstlisting}

\subsection{BlockResidual}

\begin{lstlisting}
ResidualAReal[i]   = RhsAReal[i]   - LhsAReal[i];
ResidualAImag[i]   = RhsAImag[i]   - LhsAImag[i];
ResidualPhiReal[i] = RhsPhiReal[i] - LhsPhiReal[i];
ResidualPhiImag[i] = RhsPhiImag[i] - LhsPhiImag[i];
\end{lstlisting}

\subsection{BlockAXPY}

实现
\begin{equation}
    X \leftarrow X+\alpha Y.
\end{equation}

粒子形式：
\begin{lstlisting}
ARealX[i]   += alpha * ARealY[i];
AImagX[i]   += alpha * AImagY[i];
PhiRealX[i] += alpha * PhiRealY[i];
PhiImagX[i] += alpha * PhiImagY[i];
\end{lstlisting}

\subsection{BlockDot}

实现
\begin{equation}
    \langle X,Y\rangle
    =
    \sum_i V_i
    \left[
        A_{\mathrm{re},X}\cdot A_{\mathrm{re},Y}
        +
        A_{\mathrm{im},X}\cdot A_{\mathrm{im},Y}
        +
        \phi_{\mathrm{re},X}\phi_{\mathrm{re},Y}
        +
        \phi_{\mathrm{im},X}\phi_{\mathrm{im},Y}
    \right]_i.
\end{equation}

\subsection{BiCGStabSolver}

主循环中只反复调用：
\begin{lstlisting}
AphiApply(P, V);  // V = K P
AphiApply(S, T);  // T = K S
\end{lstlisting}

并通过 block 向量操作更新：
\begin{lstlisting}
P, S, X, R
\end{lstlisting}

\section{关于调试版与最终版}

调试阶段可以直接使用 SPHinXsys 已有的
\begin{lstlisting}
LinearGradient<Real>
LinearGradient<Vecd>
\end{lstlisting}
来显式计算
\begin{equation}
    \nabla\phi,
    \qquad
    \nabla\cdot\bm A.
\end{equation}

这样方便确认符号和耦合项是否正确。

最终版建议把梯度、散度、Laplace 和耦合项都融合到 \texttt{AphiApply} 的邻居循环中，避免产生大量中间变量和多次 kernel launch。

\section{一句话总结}

A--$\phi$ 的 matrix-free 求解不是每个粒子独立求一个局部小矩阵，也不是显式组装一个 $8N\times 8N$ 的全局稀疏矩阵。它的核心是：

\begin{equation}
    \text{通过每个粒子的邻居循环直接计算 } \mathcal{K}X,
\end{equation}

然后用
\begin{equation}
    R = B-\mathcal{K}X
\end{equation}
构造残差，再用 BiCGStab 这类 Krylov 迭代方法不断更新
\begin{equation}
    X =
    \left[
    \bm A_{\mathrm{re}},
    \bm A_{\mathrm{im}},
    \phi_{\mathrm{re}},
    \phi_{\mathrm{im}}
    \right]^T,
\end{equation}
直到残差足够小。

\end{document}
